% ---
% Abstract
% ---

% resumo em inglês
\begin{resumo}[Abstract]
 \begin{otherlanguage*}{english}
This work presents the implementation and validation of the MUSE 
(Multi-Sensor State Estimator) multi-sensor state estimator on the Unitree Go2 
(EDU version) quadruped robot. The primary objective is the application of a 
robust estimation strategy that combines proprioceptive data (IMU and encoders) 
and exteroceptive data (LiDAR or camera) to ensure locomotion stability in 
complex terrains. Unlike conventional approaches, the project utilizes the DLS2 
(Dynamic Legged Systems) middleware and focuses on reducing position drift 
through dedicated slip detection modules. The theoretical framework describes 
the robot's kinematics with 12 degrees of freedom, the cascaded observer 
architecture (NLO and XKF), and sensor fusion. The estimator's initial 
validation will be performed in a simulated Gazebo environment, with telemetry 
analysis via PlotJuggler and spatial data visualization through the PCL
 (Point Cloud Library), aiming to prepare the system for real-world hardware 
 applications.

   \textbf{Keywords}: State estimation. Quadruped robot. Sensor fusion. MUSE. Unitree Go2. Slip detection.
 \end{otherlanguage*}
\end{resumo}
